% ============================================================================
%  PROJECT SUDARSHAN - Prototype Progress & Engineering Benchmark Report
%  Formatted to the supplied IEEE prototype template (spconf.sty + IEEEbib.bst)
%
%  Build:  pdflatex Sudarshan_Prototype_Report.tex   (x2 for cross-references)
%
%  EVIDENTIARY POLICY FOR THIS DOCUMENT
%  ------------------------------------
%  Every quantitative claim below is traceable to a file path and line in the
%  Sudarshan repository, or to the machine-generated corpus record
%  docs/evaluation/corpus_static_validation.json (results digest
%  1f9cd4c2...ea5f6e). Where the project's own design notes disagree with the
%  implemented code, Section 7 states the discrepancy explicitly and the code
%  is treated as authoritative.
% ============================================================================
\documentclass{article}
\usepackage{spconf,amsmath,amssymb,graphicx,hyperref,booktabs,array}

\title{PROJECT SUDARSHAN: A DETERMINISTIC, EVIDENCE-GATED ENGINE FOR MOBILE
BANKING FRAUD ANALYSIS AND VISUAL IMPERSONATION DEFENCE}

\name{Prototype Progress and Engineering Benchmark Report}
\address{Project Sudarshan --- Autonomous Mobile Banking Fraud Analysis Platform}

\begin{document}
\ninept
\maketitle

\begin{abstract}
Android banking trojans defeat commercial multi-engine scanners in two ways:
they ship a clean stub and unpack the real payload at runtime, and they stay
dormant under observation so a sandbox records nothing. Both failure modes are
scored identically by conventional engines --- as an absence of malice. This
report documents Sudarshan, a working prototype that treats an empty
observation as an unresolved question rather than an acquittal. Three
contributions are measured: a Visual Impersonation Detection Engine that
identifies bank interface clones by perceptual colour distance and normalised
view-hierarchy structure rather than by code signatures; a deterministic
four-axis Fraud Risk Score whose axes are excluded and renormalised when they
carry no data; and an Execution Assertion Matrix recording which of a sample's
own trigger preconditions a sandbox run actually reached. On a 17-sample
labelled corpus the engine attains precision and recall of 1.00 at its
deployment gate, although the numeric score ranges of the two classes overlap
and the separation is produced by evidence-gating rules rather than by the
score. That limitation, and seven others, are stated in full.
\end{abstract}

\begin{keywords}
Android banking trojan, visual impersonation, CIEDE2000, dynamic
instrumentation, evidence gating, MITRE ATT\&CK for Mobile
\end{keywords}

% ===========================================================================
\section{Introduction}
\label{sec:intro}

A mobile banking trojan does not need to be sophisticated to be effective. It
needs to look like a bank, read one SMS message, and be running when the victim
opens the real application. The Indian retail banking threat surface --- Drinik,
SOVA, and the localised Cerberus and Octo derivatives --- is dominated by
samples whose static footprint is deliberately unremarkable and whose runtime
behaviour is deliberately conditional.

This defeats the standard analysis pipeline at two specific points.

\textbf{Static analysis fails on droppers.} The manifest describes a stub. In
the labelled corpus used in Section~\ref{sec:eval}, the Hook sample declares
\emph{one} permission and zero services, and Anubis declares four. A
capability-weighted static score computed over those manifests ranks both
samples below a file manager, because the code that will actually run is not in
the file being analysed.

\textbf{Dynamic analysis fails on dormancy.} A sandbox that runs a sample for
ninety seconds and records nothing has produced an ambiguous result: either the
sample is benign, or its preconditions were never met. Anatsa waits for a
targeted banking package to enter the foreground. SOVA waits for an
accessibility grant. An SMS stealer waits for a message to arrive. None of these
occur in a sterile emulator, and a scoring model that reads the resulting
silence as a low behavioural score rewards evasion directly: the better the
sample hides, the safer it is rated.

Sudarshan is an attempt to build a pipeline in which neither failure is silent.
Its design commitments are:

\begin{enumerate}
\item \textbf{Detect the interface, not the code.} A phishing overlay that
reproduces a bank's login screen is recognisable as a clone even when its
bytecode is entirely novel. Section~\ref{sec:vide} describes the engine that
does this perceptually and structurally.
\item \textbf{Make the formula legible.} Every score is a weighted sum of named
axes with published weights and per-axis evidence strings; there is no learned
model whose verdict cannot be recomputed by hand (Section~\ref{sec:math}).
\item \textbf{Never score an absence as a zero.} An axis with no data is
removed from the formula and the remaining weights are renormalised, rather
than contributing a benign-looking zero at full weight
(Section~\ref{ssec:frs}).
\item \textbf{Record what the sandbox failed to reach.} The Execution Assertion
Matrix (Section~\ref{ssec:eam}) enumerates six fraud preconditions and reports,
per run, which were reached --- so a quiet run is labelled
\texttt{INCOMPLETE\_EXERCISE} rather than \emph{clean}.
\end{enumerate}

The prototype comprises 120 Python modules under
\texttt{shared/sudarshan\_core/}, a 2\,309-line Frida instrumentation agent, 122
React/TypeScript components, and 92 Python test modules. All measurements in
this report were taken from that tree at commit \texttt{ec7e6ff}.

% ===========================================================================
\section{Threat Model and System Architecture}
\label{sec:arch}

\subsection{Adversary model}
\label{ssec:adversary}

The adversary controls the APK entirely and is assumed to know that automated
analysis is occurring. Specifically, the adversary may: (i) ship arbitrary
encrypted payloads as assets and decrypt them into memory at runtime; (ii)
obfuscate all identifiers and construct all sensitive strings by reflection;
(iii) fingerprint the analysis environment via \texttt{SystemProperties},
\texttt{Build} fields, \texttt{/proc} inspection, and debugger checks; (iv)
gate all malicious behaviour behind preconditions that a sterile sandbox does
not satisfy; and (v) reproduce a bank's user interface pixel-for-pixel under a
package name and signing certificate of the adversary's choosing.

The adversary is \emph{not} assumed to be able to defeat in-process
instrumentation once the payload executes, nor to be able to make a cloned
banking interface visually distinct from the interface it is cloning --- the
attack requires the victim to be deceived, which bounds how far the clone can
diverge. Those two residual constraints are what Sections~\ref{sec:vide}
and~\ref{sec:dyn} exploit.

\subsection{Pipeline}
\label{ssec:pipeline}

Analysis proceeds through five stages, each of which emits structured evidence
records rather than a scalar:

\begin{enumerate}
\item \textbf{Static intake.} Dual decompilation via a MobSF container client
(\texttt{services/mobsf\_client.py}) and a native Androguard-based analyser
(\texttt{analyzers/apk\_analyzer.py}), producing the manifest permission set,
certificate block, dangerous-API inventory, string-pool entropy, and hardcoded
network indicators.
\item \textbf{Visual profiling (VIDE).} Extraction of a UI profile ---
strings, view sequence, colour palette --- from layout XML, WebView HTML,
JavaScript bundles, and asset stylesheets
(\texttt{engines/vide/pipeline.py}), then comparison against institution
baselines.
\item \textbf{Dynamic instrumentation.} Frida attachment with ART
deoptimisation, 67 distinct Java hooks, clock control, and synthetic device
state (Section~\ref{sec:dyn}).
\item \textbf{Threat correlation.} Concurrent lookups against VirusTotal
(file hash, URLs), AlienVault OTX (hash, domains), and AbuseIPDB (IPs), with a
24-hour indicator cache (\texttt{services/threat\_correlator.py}). The axis is
marked unavailable, not zero, when no API keys are provisioned.
\item \textbf{Scoring and synthesis.} The deterministic risk engine
(Section~\ref{sec:math}) followed by a retrieval-grounded Gemini narrative
(\texttt{backend/app/ai/gemini\_rag.py}, model \texttt{gemini-2.5-flash}) that
is constrained to an evidence index built from the case record.
\end{enumerate}

Runtime state is exposed to the investigator dashboard through eight endpoints
under \texttt{/runtime/} (\texttt{backend/app/routes/runtime\_api.py}), covering
health, hook inventory, live events, pipeline stage, metrics, evidence, and
diagnostics.

% ===========================================================================
\section{The Visual Impersonation Detection Engine}
\label{sec:vide}

\subsection{Problem statement}
\label{ssec:vide-problem}

A credential-harvesting overlay is a user-interface artefact. Its bytecode may
be freshly written, its strings may be constructed at runtime, and its assets
may be fetched from a C2 server after installation --- in which case every
signature, permission-count, and entropy heuristic returns nothing of interest.
What the attack cannot vary is the appearance, because a login screen that does
not look like the bank's login screen does not harvest credentials.

VIDE therefore compares the \emph{rendered identity} of a suspect application
against baselines for protected institutions, on three axes, and then applies an
entitlement test.

\subsection{Perceptual colour distance}
\label{ssec:vide-color}

Exact hexadecimal equality is the wrong comparison. A clone redrawn in a design
tool, or a palette recovered from a screenshot, lands a few units away from the
official brand value while remaining visually identical to the victim.
Conversely, raw RGB Euclidean distance is not perceptually uniform: the same
numeric distance is far more visible in some regions of the colour space than
others, and bank palettes are dominated by saturated blues and reds --- exactly
where the older $\Delta E_{76}$ and $\Delta E_{94}$ formulae misbehave.

The implementation (\texttt{engines/vide/color\_match.py}) therefore converts
sRGB to CIELAB under the D65 white point and computes the full CIEDE2000
difference. For two CIELAB colours the distance is

\begin{equation}
\Delta E_{00} = \sqrt{
\left(\frac{\Delta L'}{k_L S_L}\right)^{2} +
\left(\frac{\Delta C'}{k_C S_C}\right)^{2} +
\left(\frac{\Delta H'}{k_H S_H}\right)^{2} +
R_T \frac{\Delta C'}{k_C S_C}\frac{\Delta H'}{k_H S_H}}
\label{eq:de2000}
\end{equation}

with $k_L = k_C = k_H = 1$, and with the chroma-dependent $a^*$ correction,
the rotation term $R_T$, and the weighting functions $S_L, S_C, S_H$
implemented per CIE 142-2001.

The match strength is a graded decay rather than a hard cut, so that a
near-miss is distinguishable from an unrelated colour:

\begin{equation}
m(\Delta E_{00}) =
\begin{cases}
1 & \Delta E_{00} \le \tau_{\mathrm{id}} \\[4pt]
1 - \dfrac{\Delta E_{00} - \tau_{\mathrm{id}}}{\tau_{\max} - \tau_{\mathrm{id}}}
  & \tau_{\mathrm{id}} < \Delta E_{00} < \tau_{\max} \\[8pt]
0 & \Delta E_{00} \ge \tau_{\max}
\end{cases}
\label{eq:matchscore}
\end{equation}

with $\tau_{\mathrm{id}} = 3.5$ and $\tau_{\max} = 18.0$.

These two constants are the only tuned values in the colour axis, and both are
anchored to published perceptual quantities rather than chosen by inspection.
$\Delta E_{00} \approx 1.0$ is the just-noticeable difference under reference
conditions and $\approx 2.3$ is the threshold at which a non-expert observer
reliably reports two swatches as different; $\tau_{\mathrm{id}} = 3.5$ sits
just above that, so a palette that survived a screenshot round-trip still
matches. $\tau_{\max} = 18.0$ stays inside the same colour family, so a bank
with a merely adjacent palette cannot outrank the institution the suspect
actually reproduced.

Achromatic colours are excluded before comparison (saturation $< 0.12$, plus
near-white and near-black rejection): greys carry no brand identity because
every banking application uses them, and admitting them would let background
fills dominate the score. Palette similarity is then computed as mean coverage
of the baseline's brand colours,

\begin{equation}
S_{\mathrm{colour}} = \frac{1}{|B|}\sum_{c \in B}
\max_{s \in S} \; m\!\left(\Delta E_{00}(c, s)\right),
\label{eq:palette}
\end{equation}

where $B$ is the baseline brand palette and $S$ the suspect's. The formulation
is deliberately coverage-oriented and asymmetric: the question is whether the
suspect wears the bank's colours, not whether the two palettes are the same
size, so a suspect carrying additional unrelated colours is not penalised.

\subsection{View-hierarchy normalisation and structural similarity}
\label{ssec:vide-ast}

A clone is rarely built with the same toolkit as its target. A Capacitor or
React Native clone of a native application renders \texttt{<input>} where the
original declares \texttt{TextInputEditText}. A comparison over raw tag names
therefore fails on precisely the cases that matter.

\texttt{engines/vide/view\_ast.py} resolves this by translating every view tag
--- from Android layout XML, UIAutomator dumps, or WebView HTML --- into a
nine-symbol role vocabulary: \texttt{CONTAINER}, \texttt{TEXT}, \texttt{INPUT},
\texttt{BUTTON}, \texttt{IMAGE}, \texttt{LIST}, \texttt{WEBVIEW}, \texttt{NAV},
\texttt{OTHER}. Translation is idempotent, so an already-normalised sequence can
be re-translated safely.

The tree is retained rather than flattened. A parenthesised skeleton such as
\texttt{CONTAINER(TEXT,INPUT,INPUT,BUTTON)} encodes nesting, which a flat tag
list cannot --- without it, a settings page and a login form with equal widget
counts produce identical signatures. Structural similarity blends a
nesting-sensitive edit ratio with a role-histogram cosine:

\begin{equation}
S_{\mathrm{struct}} = 0.45 \cdot \mathrm{ratio}\!\left(\sigma_a, \sigma_b\right)
 + 0.55 \cdot \cos\!\left(\mathbf{h}_a, \mathbf{h}_b\right),
\label{eq:tree}
\end{equation}

where $\sigma$ is the skeleton string and $\mathbf{h}$ is the role histogram
weighted by discriminative power: $w_{\texttt{INPUT}} = 3.0$,
$w_{\texttt{BUTTON}} = w_{\texttt{LIST}} = w_{\texttt{NAV}} =
w_{\texttt{WEBVIEW}} = 2.0$, $w_{\texttt{IMAGE}} = 1.0$,
$w_{\texttt{TEXT}} = 0.5$, $w_{\texttt{CONTAINER}} = 0.3$,
$w_{\texttt{OTHER}} = 0.2$.

The weight ordering is the substantive design decision here, and it is a claim
about \emph{discriminative power, not importance}. Containers and loose text
appear in every screen ever built; counting them equally with input fields makes
any two screens look alike. Input fields are weighted six times a text label
because a two-field credential form is the structure the attack requires.

The engine additionally infers named structural signatures ---
\texttt{AUTH\_FORM\_VERTICAL\_PRIMARY\_CTA}, \texttt{PINPAD\_NUMERIC\_ENTRY},
\texttt{DASHBOARD\_CARD\_GRID\_QUICKACTIONS} --- from widget shape first and
text only as corroboration, so a clone that localises its labels still matches.
The PIN-pad inference requires either nine or more digit-labelled buttons, or
six or more together with PIN vocabulary, or an explicit entry prompt on a
screen that is not already a multi-field credential form; this last guard exists
because \emph{``Forgot MPIN?''} on a login screen is a hyperlink, not a keypad.

\subsection{Fusion and the entitlement predicate}
\label{ssec:vide-fusion}

The three axes are fused by \texttt{engines/vide/compare.py} into rule
\texttt{VIDE-F001}:

\begin{equation}
C_{\mathrm{VIDE}} = 0.40 \cdot S_{\mathrm{str}}
 + 0.35 \cdot S_{\mathrm{struct}}
 + 0.25 \cdot S_{\mathrm{colour}},
\label{eq:vide}
\end{equation}

where $S_{\mathrm{str}}$ is fuzzy token-set containment of the baseline's
labels, so that a reworded clone still scores.

Strings carry the largest weight because they are the most institution-specific
of the three signals --- ``YONO'' and ``Enter UPI PIN'' identify a bank in a way
that a blue palette and a two-field form do not. Structure is second because it
is toolkit-independent after normalisation but shared across the sector. Colour
is weighted lowest despite being the most perceptually salient signal, because
it is the easiest for a competent adversary to shift deliberately while
remaining convincing to a victim.

Detection requires three conditions simultaneously:

\begin{equation}
D = \big[C_{\mathrm{VIDE}} \ge 0.72\big] \wedge
    \big[S_{\mathrm{str}} \ge 0.08 \vee S_{\mathrm{struct}} \ge 0.35\big] \wedge
    \neg O,
\label{eq:vide-detect}
\end{equation}

The second conjunct is a guard against confidence accumulated entirely from the
structural axis, which every banking application shares; a finding must show
that the app reproduces \emph{this} institution's text or palette, not merely
that it is shaped like a banking app.

The third conjunct, $O$, is the entitlement predicate, and it is what makes the
rule sound. Reproducing a bank's interface is exactly what that bank's own
application does. What makes reproduction into impersonation is doing so without
being entitled to the identity. \texttt{claims\_official\_identity()} returns
true only when the suspect carries both a registered package name for the
institution \emph{and} a signing certificate on file for it. A matching package
name alone proves nothing --- anyone may build an APK declaring
\texttt{com.sbi.lotus}; it is the signature that makes the claim true.

\subsection{CH06: signer impersonation}
\label{ssec:ch06}

Rule \texttt{CH06-SIGNER-IMPERSONATION}
(\texttt{engines/vide/signer\_registry.py}) enforces the identity claim
independently of any visual finding. The registry
(\texttt{data/bank\_signer\_registry.json}) holds 12 protected Indian banking
packages: \texttt{com.sbi.lotus} and \texttt{com.sbi.lotusintouch} (YONO SBI),
\texttt{com.snapwork.hdfc} and \texttt{com.hdfc.bank},
\texttt{com.csam.icici.bank.imobile}, \texttt{com.axis.mobile},
\texttt{com.mconnect} (bob World), \texttt{com.pnb.pnbone},
\texttt{com.boi.mobile} (BOI Mobile), \texttt{com.msf.k8} (Kotak811),
\texttt{com.indusind.indusmobile}, and \texttt{com.infrasofttech.uboi} (Vyom).

The rule fires when a package in the registry is signed by a certificate not on
its allowlist. Two behaviours are worth stating precisely because they determine
how the rule behaves in the current deployment:

\emph{An empty allowlist fails closed.} Where no canonical fingerprint has been
provisioned, the identity claim cannot be verified and is therefore rejected.
This is a deliberate state documented in the registry file, not missing data.
The consequence for the present prototype is material and is restated in
Section~\ref{sec:limits}: \textbf{all 12 registry entries currently hold zero
provisioned fingerprints}, so CH06 today rejects \emph{any} APK claiming a
protected package name, and is not yet performing the certificate
\emph{comparison} the design describes.

\emph{An unextractable fingerprint is not a finding.} If no SHA-256 can be
recovered from the certificate block, the rule reports
\texttt{CH06 undetermined} rather than asserting impersonation on absent
evidence --- unless the certificate is self-evidently a debug or test keystore
(matched on subject patterns such as \texttt{CN=Android Debug}), which is
conclusive without a fingerprint, since such an APK never passed through a
bank's release pipeline.

\subsection{Deterministic escalation: the CH27 triad}
\label{ssec:ch27}

A high-confidence visual clone, signed by someone other than the bank, that
also requests the accessibility service, is the complete on-device fraud chain:
it looks like the bank, it is not the bank, and it can drive the screen.
\texttt{risk\_engine.py} escalates this combination to a floor of 95.0 and band
\texttt{Critical}:

\begin{equation}
\mathrm{CH27} =
\big[C_{\mathrm{VIDE}} > 0.85\big] \wedge
\big[\mathrm{CH06}\big] \wedge
\big[\texttt{BIND\_ACCESSIBILITY\_SERVICE}\big].
\label{eq:ch27}
\end{equation}

Each signal alone is defensible --- a skinned theme, a re-signed enterprise
build, a legitimate screen reader --- which is why \texttt{Critical} is reserved
for the conjunction. Signer impersonation alone floors at 92.0; a critical
visual cluster alone floors at 88.0.

% ===========================================================================
\section{Dynamic Instrumentation and Anti-Evasion}
\label{sec:dyn}

\subsection{Hook surface}
\label{ssec:hooks}

The instrumentation agent (\texttt{engines/frida\_hooks/banking\_trojan.js},
2\,309 lines) installs 67 distinct Java hooks across 72 registration sites,
grouped into scored and unscored evidence categories. The fraud-relevant
coverage is:

\begin{table}[htb]
\centering
\caption{Instrumentation coverage by attack function (selected).}
\label{tab:hooks}
\begin{tabular}{@{}p{2.05cm}p{5.4cm}@{}}
\toprule
\textbf{Function} & \textbf{Hooked members} \\
\midrule
Accessibility &
\texttt{AccessibilityService.onAccessibilityEvent}, \texttt{.dispatchGesture},
\texttt{AccessibilityNodeInfo.getText}, \texttt{.performAction},
\texttt{.findAccessibilityNodeInfosByText} \\
\addlinespace
SMS / OTP &
\texttt{SmsManager.sendTextMessage}, \texttt{.sendMultipartTextMessage},
\texttt{SmsMessage.getMessageBody},
\texttt{ContentResolver.query} (sms/mms/contacts), \texttt{.insert},
\texttt{.delete} \\
\addlinespace
Overlay &
\texttt{WindowManager.addView}, \texttt{.updateViewLayout},
\texttt{.removeView}, \texttt{InputMethodManager.showSoftInput} \\
\addlinespace
Payload drop &
\texttt{DexClassLoader.<init>}, \texttt{PathClassLoader.<init>},
\texttt{InMemoryDexClassLoader.<init>}, \texttt{ClassLoader.loadClass},
\texttt{System.loadLibrary} \\
\addlinespace
Persistence &
\texttt{AlarmManager.setExact}, \texttt{JobScheduler.schedule},
\texttt{DevicePolicyManager.isAdminActive}, \texttt{.lockNow} \\
\addlinespace
C2 / exfil &
\texttt{HttpURLConnection.getInputStream}, \texttt{HttpsURLConnection.connect},
\texttt{OkHttp.RealCall.execute}/\texttt{.enqueue},
\texttt{Retrofit.OkHttpCall.execute}, \texttt{Socket.connect},
\texttt{URL.openConnection} \\
\addlinespace
Evasion (unscored) &
\texttt{SystemProperties.get}, \texttt{Debug.isDebuggerConnected},
\texttt{Build.staticFields}, \texttt{BufferedReader.readLine} (anti-Frida),
\texttt{ActivityManager.getRunningTasks} \\
\bottomrule
\end{tabular}
\end{table}

Evasion hooks are deliberately placed in an \emph{unscored} category
(Section~\ref{ssec:bfci}). A sample fingerprinting the sandbox has told us
something important about itself, but it has not committed a fraud act, and
scoring it as one would inflate every verdict against every sample that checks
its environment --- including legitimate applications with anti-tamper logic.

\subsection{Time-warp clock control}
\label{ssec:timewarp}

FluBot and SharkBot variants stall behind multi-minute sleeps and
\texttt{AlarmManager}/\texttt{WorkManager} timers specifically to outlast a
sandbox window. \texttt{engines/frida\_hooks/time\_warp.js} addresses this by
applying an in-process offset to every clock the process can read:
\texttt{System.currentTimeMillis}, \texttt{SystemClock.elapsedRealtime},
\texttt{.uptimeMillis}, \texttt{.elapsedRealtimeNanos},
\texttt{java.util.Date.<init>()}, \texttt{Calendar.getInstance()}, and
\texttt{java.time.Instant.now()}. The offset is controlled at runtime through
the RPC surface \texttt{setTimeOffset}, \texttt{advanceHours},
\texttt{getTimeOffset}, \texttt{status}, and \texttt{reset}.

Three design points are worth recording, since they are what distinguish this
from the obvious alternative of moving the device clock with
\texttt{adb shell date}:

\begin{itemize}
\item \texttt{elapsedRealtime()} and \texttt{uptimeMillis()} are monotonic since
boot and do \emph{not} move when the wall clock moves. A sample that measures a
delay with them is entirely unaffected by any amount of \texttt{date}
manipulation. Hooking the accessors moves the monotonic clocks too.
\item Setting the system clock is a device-wide, externally observable change.
A sample that samples the clock twice per second sees it lurch --- which is
itself a sandbox tell. An in-process offset is visible only to the process under
analysis and leaves the rest of the device self-consistent.
\item Overloads taking an explicit timestamp are deliberately \emph{not}
shifted. \texttt{new Date(storedMillis)} is the sample restating a value it
already holds; offsetting it would corrupt the sample's arithmetic rather than
advance its timers.
\end{itemize}

Every hook is installed inside an independent guard, because the overload sets
of \texttt{Date} and \texttt{Calendar} vary across Android 7 through 15. A
missing class degrades to ``that one clock is unhooked'', never to ``the script
failed to load'', and the \texttt{status} RPC reports exactly which clocks are
under control on the current image. Unit coverage is at
\texttt{tests/unit/test\_time\_warp.py}.

\subsection{Synthetic persona seeding}
\label{ssec:persona}

An out-of-the-box emulator has no contacts, no messages, and no call history.
This is a stronger sandbox tell than any build property, because checking
whether \texttt{ContactsContract} returns zero rows costs one query and cannot
be defeated by renaming a fingerprint.

\texttt{engines/persona.py} expands operator-editable JSON templates
(\texttt{config/personas/*.json}) into concrete provider rows. Two personas ship:

\begin{table}[htb]
\centering
\caption{Shipped persona templates.}
\label{tab:persona}
\begin{tabular}{@{}lrrrr@{}}
\toprule
\textbf{Persona} & \textbf{Contacts} & \textbf{SMS} & \textbf{Calls} & \textbf{Photos} \\
\midrule
\texttt{default\_retail\_user} & 120 & 8 & 24 & 18 \\
\texttt{minimal\_privacy\_user} & 25 & 2 & 6 & 4 \\
\bottomrule
\end{tabular}
\end{table}

Names, family names and phone prefixes are drawn from India-localised pools, so
the seeded state is consistent with the device locale a targeted sample expects.
Expansion is deterministic --- seeded from the persona identifier rather than
the clock --- so the same persona yields the same device state on every run.
This is a forensic requirement, not a convenience: a re-executed run must land
on identical device state for its result to be comparable, and the module
performs no device I/O so that this property is unit-testable without an
emulator. Row counts are bounded (500 contacts, 200 SMS, 300 calls, 100 photos)
so that an operator typo cannot become a hanging insert loop.

\subsection{Runtime payload interception}
\label{ssec:dex}

Cerberus- and Octo-lineage droppers ship clean stub code and decrypt the real
payload into memory at runtime. Sudarshan hooks all three loader paths ---
\texttt{DexClassLoader.<init>}, \texttt{PathClassLoader.<init>}, and
\texttt{InMemoryDexClassLoader.<init>} --- capturing the loaded path or byte
buffer at the moment of construction, together with
\texttt{ClassLoader.loadClass} for subsequent resolution. The static side of the
same problem is handled by the concealed-payload detector, which contributes
$+60$ to the obfuscation axis (Section~\ref{ssec:stei}) and, more importantly,
arms the visibility floor described in Section~\ref{ssec:floors}.

% ===========================================================================
\section{Deterministic Risk Formulation}
\label{sec:math}

\subsection{Static Threat and Environmental Index}
\label{ssec:stei}

STEI aggregates five static axes, each normalised to $[0,100]$:

\begin{equation}
\mathrm{STEI} = 0.60\,C + 0.20\,B + 0.10\,P + 0.05\,O + 0.05\,I
\label{eq:stei}
\end{equation}

where $C$ is credential theft, $B$ banking targeting, $P$ permission risk, $O$
obfuscation, and $I$ infrastructure risk.

The 0.60 mass on $C$ is the central modelling assertion: for banking fraud
specifically, the three credential-harvesting vectors are not merely the most
common indicators, they are the mechanism by which money actually moves. $C$
decomposes as \texttt{BIND\_ACCESSIBILITY\_SERVICE} $+40$ (screen scraping and
tap injection), SMS read/receive $+35$ (OTP interception), and
\texttt{SYSTEM\_ALERT\_WINDOW} $+25$ (phishing overlay) --- the ordering
reflects that accessibility subsumes the other two capabilities in practice,
since a service that can read the screen and inject taps does not strictly need
either of the others.

$B$ is driven by matched Indian banking package names ($20$ base, $+10$ per
package) and is additionally satisfied by a VIDE visual finding at
$\min(35 + 40\,C_{\mathrm{VIDE}},\, 85)$ --- so an application that clones a
bank's interface scores as targeting that bank even when it references no
banking package name at all, which is the dropper case.

$O$ is where the model deviates most from a naive capability count, and
deliberately. A concealed executable payload contributes $+60$, dominating the
axis, because it defeats static analysis outright: the manifest describes a
stub, not the code that will run. The code comment records the empirical basis
--- Anubis, Hook and Drinik declared 4, 1 and 16 permissions respectively and
scored below a file manager before this term existed.

\subsection{Behavioural Fraud Confidence Index (v2)}
\label{ssec:bfci}

BFCI scores observed runtime behaviour over six weighted categories
(\texttt{engines/bfci\_scorer.py}):

\begin{equation}
\mathrm{BFCI}_{\mathrm{raw}} = 0.35\,A + 0.25\,S + 0.20\,W
 + 0.10\,B_{\mathrm{dyn}} + 0.05\,N + 0.05\,R
\label{eq:bfci}
\end{equation}

with $A$ accessibility, $S$ SMS/OTP, $W$ overlay window, $B_{\mathrm{dyn}}$
banking-app targeting, $N$ C2 network, and $R$ persistence.

Version 2 replaced a distinct-hook-name count with a three-tier model, because
v1 scored three unrelated accessibility events identically to a real OTP-theft
chain. Each component is now volume-aware on a logarithmic scale:

\begin{equation}
X_{\mathrm{cat}} = 100 \cdot \min\!\left(
\frac{\ln(n+1)}{\ln(\kappa_{\mathrm{cat}}+1)},\, 1\right),
\label{eq:component}
\end{equation}

where $n$ is the event count and $\kappa$ the per-category cap
($\kappa_A = 3$, $\kappa_S = 2$, $\kappa_W = 2$, $\kappa_B = 3$,
$\kappa_N = 10$, $\kappa_R = 2$). The logarithm is chosen so a single event
scores near 50 --- sensitive enough that first activity registers --- while
repeated activity still separates from it, which a linear scale at these small
caps cannot do.

A temporal bonus is applied when events from a fraud-relevant set of categories
co-occur within a 30-second window:

\begin{equation}
\mathrm{BFCI} = \min\!\left(\mathrm{BFCI}_{\mathrm{raw}} \cdot \mu,\ 100\right),
\quad \mu = 1.25 \text{ if a chain fires.}
\label{eq:seq}
\end{equation}

Four chains are defined: \texttt{OTP\_THEFT\_CHAIN}
($A \wedge S \wedge N$), \texttt{OVERLAY\_BANKING\_CHAIN}
($W \wedge B$), \texttt{ACCOUNT\_TAKEOVER\_CHAIN} ($A \wedge W \wedge S$), and
\texttt{DROPPER\_CHAIN} ($R \wedge N$). Ordering in time is what distinguishes
a fraud sequence from coincident API use.

Six categories are collected as evidence but never scored:
\texttt{dangerous\_apis}, \texttt{files\_accessed}, \texttt{anti\_analysis},
\texttt{device\_fingerprint}, \texttt{app\_telemetry}, \texttt{notification}.
The rationale is quantitative: given the caps above, a \emph{single} event
scores 50--63 for its component, so a scored category that also matches
ordinary application behaviour is not a weak signal --- it is a constant, and it
inflates every verdict equally. The module asserts at import that the scored and
unscored sets are disjoint.

\subsection{Fraud Risk Score and axis renormalisation}
\label{ssec:frs}

The implemented top-level score (\texttt{engines/risk\_engine.py}) is a
four-axis weighted mean over the axes that carry data:

\begin{equation}
\mathrm{FRS} = M_{\mathrm{AI}} \cdot
\frac{\sum_{k \in \mathcal{L}} w_k v_k}{\sum_{k \in \mathcal{L}} w_k},
\label{eq:frs}
\end{equation}

with nominal weights $w_{\mathrm{STEI}} = 0.25$, $w_{\mathrm{dyn}} = 0.35$,
$w_{\mathrm{corr}} = 0.20$, $w_{\mathrm{bank}} = 0.20$, and $\mathcal{L}$ the
set of \emph{live} axes. $M_{\mathrm{AI}}$ is the classification confidence
multiplier, clamped to $[0.5, 1.5]$.

Equation~\ref{eq:frs} is the Absence-of-Evidence guard, and the renormalising
denominator is the whole of it. An axis with no data is removed from
$\mathcal{L}$ rather than contributing $w_k \cdot 0$. Two axes are gated:

\textbf{Correlation} is live only when a threat-intelligence provider actually
answered. Previously an unconfigured provider contributed $0.20 \times 0.0$,
pinning 20--25\% of every score at zero in any deployment without API keys.
Measured on Teabot, that cost 10.3 FRS points --- 30.8 with correlation forced
to zero, against 41.1 with the axis correctly excluded.

\textbf{Dynamic} is live only when the run was \emph{conclusive}, defined as
$\mathrm{BFCI} \ge 20$, or at least one hook-derived behavioural event, or at
least one non-harness evidence record. Three things deliberately do not qualify:
screenshot and harness records (which prove the sandbox acted, not that the
sample did), anti-analysis events alone (evasion is resistance, not behaviour),
and a run in which the process started but never rendered an Activity or window.
The measured motivation is recorded in the source: Cerberus scored 38.74
(Suspicious) static-only; after a 30-second run that captured zero API calls and
zero network events, the dynamic axis took 0.437 of the renormalised weight at a
value of 5.0 and pulled the same binary to 23.98 --- \emph{Safe}. The better the
evasion, the safer the engine rated it.

This is not a rule that dynamic evidence may only raise a score. A run with real
coverage that observes benign behaviour is legitimate evidence and still lowers
the verdict. The distinction is solely whether the sandbox observed anything to
reason about. The exclusion reason is preserved and surfaced
(\texttt{DYNAMIC\_UNAVAILABLE}, \texttt{NO\_UI\_RENDERED},
\texttt{EVASION\_ONLY}, \texttt{INSTRUMENTATION\_FAILED},
\texttt{NO\_BEHAVIOR\_OBSERVED}), because ``nothing bad happened'' and ``we
never got to look'' are opposite claims that must not render identically.

\subsection{Verdict floors}
\label{ssec:floors}

Four floors refuse a \texttt{Safe} certification without inventing a score. In
every case the numeric FRS is left untouched and an explanatory evidence line is
emitted; only the band changes.

\begin{enumerate}
\item \textbf{Visibility floor.} Fires when a concealed payload was detected and
the sandbox did not observe substantive payload behaviour. The threshold is
deliberately a \emph{substantive} signal rather than any signal: Anubis produced
one incidental API call --- enough to mark a run conclusive --- while BFCI
stayed at 0.0 because the second stage never deployed, and a banking trojan
scored 8.72 \emph{Safe}. Octo scored BFCI 4.8 (its network component maxed at 96
but network carries weight 0.05) and that 4.8 was enough to lift the floor and
certify a heavily obfuscated sample with 31 permissions as \emph{Safe} at 29.49.
The bar is therefore $\mathrm{BFCI} \ge 20$.
\item \textbf{Static evidence floor.} Fires when $\mathrm{STEI} \ge 50$ and the
sandbox produced no fraud-relevant signal at all. A 90-second window that does
not trigger SMS interception is not evidence that the interception is absent;
the permissions remain declared.
\item \textbf{Evasion floor.} Fires when anti-analysis events occurred and
nothing else did. Because evasion carries no BFCI weight, the only observable
act on such a run contributes exactly 0.0 --- so without this floor, active
counter-analysis working as designed produces a \emph{Safe} verdict.
\item \textbf{Incomplete-exercise floor.} Fires per Section~\ref{ssec:eam}.
\end{enumerate}

\subsection{The Execution Assertion Matrix}
\label{ssec:eam}

\texttt{engines/execution\_assertions.py} records, per run, which of six fraud
preconditions were reached: targeted banking app opened, OTP/bank SMS delivered,
accessibility service granted, overlay window drawn, contacts provider read,
call log read. Each is derived from observed telemetry by regular expressions
over Android framework identifiers --- never from a curated bank list. The
banking-app assertion is evaluated against the packages \emph{this sample}
references, resolved from its own manifest and DEX constants, because passing a
curated list would encode an assumption about which institutions matter.

A run is \texttt{INCOMPLETE\_EXERCISE} when the sandbox ran, observed zero
threat events, and reached zero assertions:

\begin{equation}
\mathrm{IE} = R_{\mathrm{dyn}} \wedge (E = 0) \wedge
\Big(\textstyle\bigwedge_{i=1}^{6} \neg a_i\Big).
\label{eq:ie}
\end{equation}

The \texttt{dynamic\_ran} conjunct matters: a case with no sandbox run is
already reported as \emph{dynamic analysis unavailable}, and labelling it an
incomplete \emph{exercise} would falsely claim we tried and the sample stayed
quiet. When \texttt{IE} holds, reported confidence is multiplied by 0.50, each
unreached assertion is emitted with a concrete remediation instruction, and a
\texttt{Safe} band is floored to \texttt{Suspicious}.

This is the mechanism that converts the industry's most common false negative
--- a clean-looking report from a run that never happened --- into an actionable
statement about the harness.

% ===========================================================================
\section{Evaluation}
\label{sec:eval}

\subsection{Labelled corpus}
\label{ssec:corpus}

The primary evaluation is a machine-generated record over a 17-sample labelled
corpus (\texttt{docs/evaluation/corpus\_static\_validation.json}, git commit
\texttt{ce30610}, Androguard 4.1.4, results digest
\texttt{1f9cd4c2\ldots ea5f6e}). The run configuration is deliberately
adversarial to Sudarshan's own strengths: \textbf{dynamic analysis not run,
correlation unavailable, VIDE not run}. Only the static and banking-impact axes
were live, at renormalised weights 0.556 and 0.444.

The decision rule is the product's own deployment gate rather than a threshold
chosen for the paper: \texttt{Safe} yields \emph{``approved for deployment''}
and any higher band yields \emph{``quarantine''}, so the band boundary is what a
consumer acts on. Deliberately-insecure applications (\texttt{InsecureBankv2})
are counted as benign, since insecure is not malicious.

\begin{table}[htb]
\centering
\setlength{\tabcolsep}{4pt}
\caption{Corpus results. Static-only configuration; dynamic, correlation and
VIDE axes all excluded.}
\label{tab:corpus}
\begin{tabular}{@{}llrrl@{}}
\toprule
\textbf{Sample} & \textbf{Class} & \textbf{STEI} & \textbf{FRS} & \textbf{Band} \\
\midrule
Teabot            & malware & 76.59 & 70.26 & High Risk \\
FluBot            & malware & 61.20 & 60.00 & Suspicious \\
Octo              & malware & 57.20 & 57.33 & Suspicious \\
SharkBot          & malware & 52.98 & 53.45 & Suspicious \\
Cerberus          & malware & 53.73 & 38.74 & Suspicious \\
Drinik            & malware & 22.82 & 21.57 & Suspicious$^{\dagger}$ \\
Hook              & malware & 11.57 & 15.32 & Suspicious$^{\dagger}$ \\
Anubis            & malware &  9.20 & 14.00 & Suspicious$^{\dagger}$ \\
\midrule
InsecureBankv2    & benign  & 30.71 & 25.95 & Safe \\
VLC               & benign  & 25.42 & 23.01 & Safe \\
NewPipe           & benign  & 24.64 & 22.58 & Safe \\
Amaze File Mgr    & benign  &  9.70 & 14.28 & Safe \\
KeePassDX         & benign  &  7.53 & 13.07 & Safe \\
UnCrackable 1--4  & benign  & $\le$3.3 & $\le$10.7 & Safe \\
\bottomrule
\end{tabular}
\\[2pt]
{\footnotesize $^{\dagger}$ Band produced by the visibility floor
(Section~\ref{ssec:floors}), not by the numeric score.}
\end{table}

At the deployment gate the result is $\mathrm{TP}=8$, $\mathrm{FP}=0$,
$\mathrm{FN}=0$, $\mathrm{TN}=9$: precision 1.00, recall 1.00, $F_1$ 1.00,
specificity 1.00.

\textbf{This number must be read with its stated caveat, which the record
itself publishes.} The malware FRS range is $[14.00, 70.26]$ and the benign
range is $[9.16, 25.95]$. \emph{These ranges overlap.} Anubis (14.00) scores
below VLC (23.01) and well below InsecureBankv2 (25.95). The corpus record's own
\texttt{score\_separable} field reads \texttt{false}, with
\texttt{separation\_source: "risk\_band floors, not numeric FRS"}. Three of the
eight true positives --- Anubis, Drinik, Hook --- are caught \emph{only} by the
visibility floor, which fired on exactly those three and on no benign sample.
Concealment detection was the discriminator: 6 of 8 malware carried a concealed
payload against 0 of 9 benign samples.

The honest statement of this result is therefore: \emph{the evidence-gating
rules classify this corpus perfectly; the underlying continuous score does not.}
Under the stricter secondary rule (band $\ge$ High Risk) recall collapses to
0.125 --- only Teabot clears it. A production deployment would run at the
Suspicious gate and accept the resulting triage volume.

Family attribution is the weakest subsystem measured: 4 of 8 malware samples
returned \texttt{Unknown}, and of the four that resolved, three were labelled
\emph{Anubis} (FluBot, Octo, Teabot) and one \emph{Hydra} (SharkBot) --- pattern
matches on shared lineage rather than correct family names.

\subsection{Third-party cross-validation: Teabot}
\label{ssec:teabot}

The sample \texttt{com.lbugkgtau.iuzdlpqzr}, SHA-256
\texttt{554e71b4\ldots c47e1ef}, was analysed independently by a
VirusTotal/Zenbox dynamic sandbox and by Dr.Web vxCube. Binary identity is
confirmed across all three: the SHA-256 matches exactly between Sudarshan and
vxCube, and SHA-1/MD5 match between Zenbox and vxCube.

\begin{table}[htb]
\centering
\caption{Cross-tool agreement, sample \texttt{554e71b4}.}
\label{tab:xtool}
\begin{tabular}{@{}p{2.8cm}p{4.6cm}@{}}
\toprule
\textbf{Dimension} & \textbf{Result} \\
\midrule
Permission extraction & 14 / 14 of the Zenbox-published set reproduced;
Sudarshan additionally surfaced \texttt{RECEIVE\_BOOT\_COMPLETED},
\texttt{REQUEST\_PASSWORD\_COMPLEXITY}, \texttt{REQUEST\_IGNORE\_BATTERY\_-
OPTIMIZATIONS}. \textbf{100\% overlap.} \\
\addlinespace
Obfuscation technique & 3 / 3 of vxCube's obfuscation indicators (method-name
obfuscation, very long method strings, reflection) corroborated by Sudarshan's
OB axis at 93.9. \textbf{100\% agreement.} \\
\addlinespace
Dynamic behaviour & Zenbox observed accessibility interaction, SMS
interception, DEX drop (\texttt{mimlwef.dex}) and live C2. Sudarshan's own run
reached 0/6 assertions and captured 0 hooks. \textbf{Divergent runtime depth.} \\
\addlinespace
C2 indicator & Zenbox: \texttt{185.215.113.31:85} with
\texttt{/api/botupdate}, \texttt{/api/getbotinjects},
\texttt{/api/getkeyloggers}. Sudarshan: a different, malformed value.
\textbf{No overlap.} \\
\addlinespace
Launcher-icon removal & Confirmed by both Zenbox and vxCube; \textbf{not
itemised by Sudarshan} --- a coverage gap. \\
\bottomrule
\end{tabular}
\end{table}

The comparison is instructive precisely because Sudarshan's own dynamic run
failed. Zenbox reached a materially deeper execution state and observed the
behaviours Sudarshan only inferred. Sudarshan's Execution Assertion Matrix
recorded 0/6 preconditions and labelled the run accordingly, rather than
reporting a low behavioural score --- which is the designed behaviour, and the
reason the sample was still escalated.

Numeric verdicts across the three tools (Sudarshan FRS 46.3; vxCube ``Zenbox
Android Verdict'' 52/100 carrying a \emph{Non Malicious} label) are
\textbf{not comparable}: they are different formulae over different evidence,
and no accuracy figure is derived from their relationship anywhere in this
report.

\subsection{Arithmetic verification of the renormalisation}
\label{ssec:arith}

The two independently produced Teabot scores validate
Equation~\ref{eq:frs} numerically. In the case-report run, the dynamic axis was
excluded and correlation was live at 18.67 with banking impact 36.00:

\begin{equation}
\frac{0.25(76.59) + 0.20(18.67) + 0.20(36.00)}{0.65} = 46.28.
\label{eq:teabot1}
\end{equation}

In the corpus run, correlation was \emph{also} excluded:

\begin{equation}
\frac{0.25(76.59) + 0.20(62.35)}{0.45} = 70.26.
\label{eq:teabot2}
\end{equation}

Both reproduce the published values exactly, from the same STEI of 76.59.

This also exposes a genuine reporting defect worth stating plainly: the case
report renders its working as
$0.25 \times 76.59 + 0.35 \times 0.00 + 0.20 \times 18.67 + 0.20 \times 36.00
= 46.28$, which does not evaluate to 46.28 --- it evaluates to 30.08. The
\emph{engine} performed the renormalisation correctly; the report template
prints the nominal weights without the denominator. The arithmetic shown to an
analyst must match the arithmetic performed.

The comparison of Equations~\ref{eq:teabot1} and~\ref{eq:teabot2} is the
clearest available demonstration of what the guard is worth: the same binary,
with the same static evidence, moves from 46.28 (\emph{Suspicious}) to 70.26
(\emph{High Risk}) purely by declining to average an empty axis into the result.

% ===========================================================================
\section{Adversarial Review: Discrepancies and Limitations}
\label{sec:limits}

The following were identified by reviewing the project's own design notes
against the implemented code. In each case the code is authoritative and the
design note is wrong. They are reported here rather than omitted because a
prototype report that overstates its own mechanisms is not usable as an
engineering baseline.

\textbf{1. The published two-axis FRS formula is not what runs.} Design notes
state $\mathrm{FRS}_{\mathrm{full}} = (0.65\,\mathrm{STEI} +
0.35\,\mathrm{BFCI}) \cdot M_{\mathrm{AI}}$ with a static fallback of
$1.0 \times \mathrm{STEI}$. The implementation is the four-axis renormalised
mean of Equation~\ref{eq:frs}, and the static-only case renormalises STEI to
0.556 and banking impact to 0.444 --- not to 1.0. This is verifiable in the
corpus record's \texttt{axes\_used} field on every sample. The implemented form
is the stronger design; the notes should be corrected to match it.

\textbf{2. Time Warp is a clock offset, not a $100\times$ accelerator.} Design
notes describe hooking \texttt{Thread.sleep}, \texttt{SystemClock.sleep} and
\texttt{Handler.postDelayed} with a $100\times$ multiplier. No such hook exists:
\texttt{time\_warp.js} hooks clock \emph{accessors} and applies an additive
offset. A sample that calls \texttt{Thread.sleep(300000)} directly still blocks
for five real minutes. The offset defeats \emph{elapsed-time comparisons}, which
is the more common dormancy pattern, but the gap is real and sleep interception
is genuine future work.

\textbf{3. The colour threshold is not $\Delta E < 10$.} The implemented band is
a graded decay from $\tau_{\mathrm{id}} = 3.5$ to $\tau_{\max} = 18.0$ over
CIEDE2000, not a hard cut at 10 over an unspecified $\Delta E$ variant
(Equation~\ref{eq:matchscore}).

\textbf{4. The MITRE identifiers in the design notes are deprecated.} Notes cite
T1412 (SMS capture), T1444 (masquerade), T1411 (input prompt) --- all retired
from ATT\&CK for Mobile in v9. The implemented mapper
(\texttt{engines/mitre\_mapper.py}) uses current identifiers and emits a valid
Navigator v4.5 layer against ATT\&CK v14: T1417.001 (Input Capture: GUI),
T1417.002 (Input Injection), T1582 (SMS Control), T1636.004 (Protected User
Data: SMS), T1407 (Download New Code at Runtime), T1626.001 (Device Admin),
T1437.001 (Web Protocols), T1634 (Credentials from Password Store), T1497.001
(Sandbox Evasion: System Checks). Two entries carry acknowledged
in-code caveats: \texttt{WindowManager.addView} is mapped to T1416, which is a
defensible but not exact fit for overlay abuse, and \texttt{lockNow} is filed
under device-admin abuse rather than a ransomware technique.

\textbf{5. CH06 is fail-closed but not yet fingerprint-verifying.} All 12
registry entries hold zero provisioned canonical signers. The rule therefore
rejects any APK claiming a protected package name, which is correct fail-closed
behaviour but is not the certificate comparison the design describes.
Provisioning the 12 release fingerprints is the single highest-value
outstanding task, and it is a data task, not an engineering one.

\textbf{6. VIDE baselines are three lab fixtures.} \texttt{data/ui\_baselines/}
contains \texttt{demo\_sbi\_yono}, \texttt{demo\_hdfc\_mobile} and
\texttt{demo\_icici\_imobile}, each versioned \texttt{1.0.0-hackathon} with
roughly ten strings, a seven-node view sequence, and a three-colour palette.
They are sufficient to exercise the engine end-to-end and to demonstrate the
mathematics, and they are not a production baseline corpus. Notably BOI is
present in the \emph{signer} registry but has no \emph{UI} baseline. VIDE was
not exercised in the corpus run of Section~\ref{ssec:corpus}, so no measured
VIDE detection rate is claimed anywhere in this report.

\textbf{7. Score separability is unproven.} Restated from
Section~\ref{ssec:corpus} because it is the most consequential limitation: the
perfect corpus result rests on evidence-gating floors, not on the continuous
score, and the class score ranges overlap. Corpus size ($n = 17$) is also small
enough that the confidence interval around a 1.00 precision estimate is wide.

\textbf{8. Published latency figures are unverified.} \texttt{BENCHMARKS.md}
carries an explicit warning that its latency, RAM and CPU figures were not
re-measured in the 2026-08-11 audit. No performance figure from that document is
cited as verified in this report.

\subsection{What the review confirms}
\label{ssec:confirms}

Set against those eight items, the following are implemented as described and
were verified directly in code and in machine-generated output: the five-axis
STEI and its weights; the six-category BFCI with logarithmic volume scaling and
temporal chain detection; axis exclusion and renormalisation with preserved
exclusion reasons; all four verdict floors; the six-assertion Execution
Assertion Matrix with per-assertion remediation text; CIEDE2000 colour distance;
role-normalised view-hierarchy similarity; the 67-hook instrumentation surface
including all three DEX loader paths; deterministic persona expansion; the
12-entry signer registry with fail-closed semantics; and the three-provider
threat correlator with cache and unavailability signalling.

% ===========================================================================
\section{Regulatory Alignment}
\label{sec:reg}

The engine emits banded actions mapped to Indian regulatory expectations
(\texttt{\_get\_recommended\_action}). A \texttt{Critical} verdict with a named
family returns an immediate-block action citing customer advisory obligations
under RBI CPG-2022; OTP interception combined with confirmed Indian bank
targeting contributes an explicit RBI MDS-2021 regulatory-risk term to the
banking-impact axis. Case records carry a chain-of-custody block and an
integrity hash, which is the property a CERT-In incident submission requires:
not that the verdict be correct, but that it be reproducible from preserved
evidence. The determinism of the scoring path is what makes this possible ---
every score in this report can be recomputed by hand from its published axes.

% ===========================================================================
\section{Conclusion and Roadmap}
\label{sec:conclusion}

Sudarshan is a working prototype whose distinguishing contribution is
methodological rather than merely functional. Commercial multi-engine scanners
converge with it on static extraction --- 100\% permission overlap with Zenbox
and 100\% obfuscation-technique agreement with vxCube on the cross-validated
sample --- and exceed it on runtime depth when their sandbox reaches a deeper
execution state. Where Sudarshan differs is in what it does with an empty
result. An axis with no data is excluded and the remaining weights renormalised;
a sandbox run that reached no trigger is labelled \texttt{INCOMPLETE\_EXERCISE}
with per-precondition remediation; a concealed payload that never deployed
floors the verdict rather than certifying it. On the labelled corpus these rules
are what separate the classes --- the continuous score alone does not.

The prototype is honest about its incompleteness. Eight discrepancies between
design intent and implementation are enumerated in Section~\ref{sec:limits},
including two --- unprovisioned signer fingerprints and three lab-fixture UI
baselines --- that gate the VIDE engine's real-world detection rate, which is
consequently not claimed.

Ordered by value delivered per unit of remaining work, the roadmap is:
(i) provision the 12 canonical bank signing fingerprints, converting CH06 from
fail-closed rejection to true certificate verification; (ii) build production UI
baselines from official Play Store releases, including BOI, and measure a real
VIDE detection rate against held-out clones; (iii) add
\texttt{Thread.sleep}/\texttt{postDelayed} interception to close the dormancy
gap identified in item 2; (iv) expand the labelled corpus and re-measure with
dynamic and VIDE axes live, to establish whether the continuous score becomes
separable once the axes that were excluded are actually populated; and (v) fix
the report template so the arithmetic shown to an analyst matches the arithmetic
the engine performs.

% ===========================================================================
\begin{thebibliography}{9}

\bibitem{cie}
CIE Technical Committee 1-47,
``Improvement to Industrial Colour-Difference Evaluation,''
CIE Publication 142-2001, Vienna, 2001.

\bibitem{sharma}
G. Sharma, W. Wu, and E. N. Dalal,
``The CIEDE2000 color-difference formula: Implementation notes, supplementary
test data, and mathematical observations,''
\textit{Color Research and Application}, vol. 30, no. 1, pp. 21--30, 2005.

\bibitem{attack}
MITRE Corporation,
``ATT\&CK for Mobile, Version 14,''
\url{https://attack.mitre.org/matrices/mobile/}, 2023.

\bibitem{threatfabric}
ThreatFabric,
``Mobile Banking Malware Landscape,'' Amsterdam, 2024.

\bibitem{frida}
O. A. V. Ravn\aa s, ``Frida: Dynamic instrumentation toolkit,''
\url{https://frida.re}.

\bibitem{androguard}
A. Desnos et al., ``Androguard: Reverse engineering and analysis of Android
applications,'' v4.1.4.

\bibitem{mobsf}
A. Abraham et al., ``Mobile Security Framework (MobSF),''
OWASP Mobile Security Testing Guide project.

\bibitem{rbi}
Reserve Bank of India,
``Master Direction on Digital Payment Security Controls,'' 2021.

\bibitem{certin}
Indian Computer Emergency Response Team (CERT-In),
``Directions under sub-section (6) of section 70B of the Information Technology
Act, 2000,'' 2022.

\end{thebibliography}

\end{document}
